\section{Principios arquitectónicos del App Builder}

\subsection{Arquitectura básica}

La arquitectura central del app builder normalmente se basa en la tecnología WebContainer: ejecuta un entorno de desarrollo completo dentro del navegador, con soporte para edición de código, compilación y vista previa.

\subsection{El motor de v0 en detalle}

Gaspar Garcia, Head of AI Research de Vercel v0, describió en detalle la arquitectura interna de v0:

\subsubsection{Flujo de cinco pasos}

\begin{enumerate}
    \item \textbf{Intent Understanding}: el procesamiento de lenguaje natural extrae los requisitos estructurados: componentes, diseño de la interfaz, interacciones, flujo de datos
    \item \textbf{Context Assembly}: recopila la información relevante: tokens del sistema de diseño, componentes existentes, esquema de la API, requisitos de accesibilidad
    \item \textbf{Code Generation}: el LLM genera componentes de React con calidad de producción (tipos de TypeScript, diseño responsivo, mejores prácticas)
    \item \textbf{Validation \& Testing}: verificaciones automatizadas del cumplimiento de accesibilidad, el comportamiento responsivo y la calidad del código
    \item \textbf{Human-in-the-Loop}: los ingenieros revisan y optimizan, y la retroalimentación se usa para entrenar el modelo
\end{enumerate}

\begin{importantbox}{El Agent es un Workflow}
Gaspar enfatizó una idea central de diseño: un Agent es, en esencia, un workflow (flujo de trabajo). La ``inteligencia'' de v0 no proviene de una sola llamada a un modelo grande, sino de un flujo de trabajo de varios pasos cuidadosamente diseñado. Cada paso tiene entradas, salidas y criterios de validación bien definidos.
\end{importantbox}

\subsection{Limitaciones del modelo de procesamiento}

\subsubsection{Stream Manipulation}

v0 creó un subsistema que puede interceptar y modificar el flujo de salida del LLM antes de que llegue al usuario. Esto resuelve el problema de que ``a veces ajustar solo el prompt no basta'': la corrección se hace directamente a nivel de la salida.

\subsubsection{Entrenamiento y ajuste fino}

Los laboratorios de modelos de frontera no crean canalizaciones de generación de código especialmente diseñadas para aplicaciones web de Next.js 16. v0 construyó su propia composite model family mediante canalizaciones propias de entrenamiento y fine-tuning, para asegurar que el código generado cumpla con las especificaciones más recientes del framework.

\begin{knowledgebox}{Composite Model Family}
v0 no depende de un solo modelo de propósito general, sino que usa un conjunto de modelos ajustados de forma específica y combinados entre sí (composite model family). Distintos modelos se encargan de distintas subtareas, y su combinación logra un resultado conjunto superior al de un solo modelo. Ver el blog de Vercel: \url{https://vercel.com/blog/v0-composite-model-family}
\end{knowledgebox}

\subsection{Resumen de la sección}

El núcleo de los App Builders modernos es un workflow de varios pasos cuidadosamente diseñado, en lugar de una sola llamada al LLM. El stream manipulation y la composite model family son los mecanismos de ingeniería clave para resolver las limitaciones de los modelos de propósito general.
